\begin{table}[t]
\centering
\scriptsize
\setlength{\tabcolsep}{4pt}
\begin{tabular}{@{}P{2.8cm} P{4.8cm} P{5.0cm}@{}}
\hline
Stress axis & Supplementary perturbation used here & Literature grounding \\
\hline
Cardiac-like pulsation\\\cite{BarangerAdaptiveSVD2018} & Reference: 1.0x cardiac amplitude at 1.30 Hz; Moderate: 2.0x cardiac amplitude at 1.30 Hz; Hard: 3.0x cardiac amplitude at 1.30 Hz & Stress axis motivated by neonatal and open-skull reports of cardiac-like tissue pulsation broadening the low-frequency tissue spectrum. \\
Short ensemble\\\cite{Imbault2017,Soloukey2020FUSAwake,MobileFUS2025} & Reference: 64 frames (42.7 ms); Moderate: 48 frames (32.0 ms); Hard: 32 frames (21.3 ms) & Stress axis motivated by clinically constrained acquisitions in which shorter ensembles are used to limit motion corruption or workflow burden. \\
\hline
\end{tabular}
\caption{Literature-grounded stress axes used in the held-out paper-tier SIMUS stress confirmation. The paper-tier confirmation fixes attention to the residualizer families that dominated the preliminary reduced-grid sweep (RPCA and adaptive-global SVD), then reruns those clinically motivated perturbations on the held-out paper-tier mobile setting. These perturbations are intended to approximate clinically relevant nuisance pressure, not to claim one-to-one replay of any single in vivo acquisition.}
\label{tab:simus_clinical_stress_frontier_provenance}
\end{table}
